%% 第6章 结论与展望

\chapter{结论与展望}

\section{研究结论}

随着大语言模型能力持续提升，智能体（AI Agent）应用在企业流程自动化、
内容生成与知识增强等场景中迅速普及。然而相对于 Python 生态中已较为成熟的
智能体编排框架，Java 生态长期缺乏原生的可视化 Agent 编排平台，开发者难以在企业
服务现有技术栈下快速构建符合业务规范的智能体应用。

针对上述问题，本文设计并实现了一套\textbf{基于 Spring Boot 的 AI Agent 智能体编排系统}，
围绕 Agent 配置、流程执行、人工审核、知识增强与访问安全五个方面完成系统设计与
实现，主要研究成果如下：

\textbf{（1）实现了 Agent 管理与可视化编排能力。}
系统通过可视化画布支持节点拖拽、参数配置、边连接与流程保存，
并对 Agent 生命周期提供创建、更新、发布、版本查询与回滚机制，
使工作流定义具备可配置、可追溯、可复用的工程化特征。

\textbf{（2）构建了事件驱动的 DAG 工作流调度执行机制。}
系统基于拓扑排序实现了“就绪节点驱动”的异步调度策略，支持无依赖节点的自动并行
执行、条件分支与路径剪枝，并通过执行聚合根统一维护节点状态、上下文数据和执行推进逻辑，
保证了工作流执行过程的可控性与可维护性。

\textbf{（3）实现了人工检查点与工作流恢复机制。}
系统支持节点执行前与执行后两种审核触发阶段，审核人可查看执行上下文并对中间结果
进行修改后放行，也可拒绝并终止流程。结合已审批节点跳过机制、分布式锁与版本号控制，
保障了多人并发审批场景下的状态一致性与审计可追踪性。

\textbf{（4）实现了知识库与检索增强链路。}
系统支持知识数据集管理、文档上传、异步解析、文本分块与向量化入库，
并基于向量数据库提供语义检索能力。在工作流运行过程中，检索结果可作为上下文注入
 LLM 节点，从而提升模型回答与业务知识之间的相关性。

\textbf{（5）实现了用户认证与系统安全支撑能力。}
系统完成了邮箱注册、登录、验证码校验、密码找回等基础身份能力，并通过密码哈希加密、
验证码发送频率限制、登录失败限流与令牌认证等机制，为 Agent 管理、知识库管理和
工作流执行等功能提供统一的访问控制基础。

综合来看，本文实现的系统已打通“可视化编排 → 工作流执行 → 人工审核 →
知识增强 → 安全访问”的核心闭环，为 Java 生态下 AI Agent 系统的落地提供了可复用的
实现思路。经测试验证，系统通过 42 个单元测试用例（核心领域层行覆盖率 78\%、
分支覆盖率 71\%）和 35 个功能测试用例的全面验证，主要管理类接口 P99 延迟低于 200\,ms，
SSE 实时推流 P95 延迟为 298\,ms，具备企业级 MVP 演示与业务验证能力。

\section{研究不足}

受时间与精力限制，本文工作仍存在以下不足：

\textbf{（1）Agent 编排能力仍以单工作流为主。}
当前系统已支持可视化编排与版本管理，但尚未支持复杂子流程复用、组合模板沉淀
与跨 Agent 协同编排，复杂业务场景的编排表达力有限。

\textbf{（2）知识增强链路较为基础。}
当前文档分块采用固定策略，未根据文档结构与领域特征进行自适应；检索阶段
未引入重排模型，且对命中片段的解释性支持不足。

\textbf{（3）人工审核交互体验有待提升。}
当前审批页面以结构化文本和 JSON 形式展示节点输入输出，对非技术用户不够友好，
缺乏富文本、表格、图片等多类型内容的差异化展示能力。

\textbf{（4）性能与高可用验证范围有限。}
现阶段测试聚焦于核心链路可用性，尚未在更高并发规模、服务重启恢复
与分布式部署等场景下进行充分验证。

\section{未来工作展望}

针对上述不足，后续可在以下方向继续研究和优化：

\textbf{（1）增强 Agent 编排的复用与组合能力。}
引入子工作流、流程模板与并行聚合节点等机制，提升复杂业务场景下的编排表达力。

\textbf{（2）完善知识检索增强效果。}
探索自适应分块策略、召回重排模型与知识来源可解释展示，提升知识命中质量
与生成结果稳定性。

\textbf{（3）优化人工审核交互与通知机制。}
结合字段级渲染、差异对比视图与邮件、即时消息通知机制，提升审核节点在
企业协作场景中的可用性与响应效率。

\textbf{（4）加强系统性能与高可用能力。}
围绕更大规模并发调度、SSE 长连接稳定性、执行断点恢复与分布式部署进行
工程优化，进一步提升系统在生产环境中的稳定性。

\textbf{（5）拓展多智能体协作与工具生态。}
在现有单 Agent 编排能力基础上，探索多智能体协同机制、更丰富的工具接入方式
以及与 Spring AI 新特性的深度整合，扩展系统的适用范围与协作能力。
